% from ega0-1-fr.tex
\oldpage[0\textsubscript{I}]{69}

de définition de $B$, autrement dit qu'il existe $k>0$ tel que
$\mathfrak n^k\subset\mathfrak mB+\mathfrak b\subset\mathfrak n$~; par suite,
pour tout $h>0$
\[
  \mathfrak n^{hk}M
  \subset(\mathfrak mB+\mathfrak b)^hM
  =\mathfrak m^hM
  \subset\mathfrak n^hM,
\]
ce qui prouve que les topologies $\mathfrak m$-préadique et
$\mathfrak n$-préadique sur $M$ sont les mêmes~; la seconde est d'ailleurs
séparée en vertu de (7.3.5).

\begin{corollary}[7.4.3]
\label{0.7.4.3-fr}
Sous les hypothèses de (7.4.2), si en outre $A$ est noethérien et complet pour
la topologie $\mathfrak m$-préadique, $M$ est un $A$-module de type fini.
\end{corollary}

En effet, $M$ est alors séparé pour la topologie $\mathfrak m$-préadique, et
notre assertion résulte de la remarque faite dans (7.4.1).

\begin{env}[7.4.4]
\label{0.7.4.4-fr}
Le cas d'application de (7.4.2), qui est le plus important, est celui où $B$
lui-même est un $A$-\emph{module quasi-fini}, ce qui revient à dire que
$B/\mathfrak mB$ est une \emph{algèbre de rang fini sur}
$k=A/\mathfrak m$~; cette condition peut d'ailleurs se décomposer en la
conjonction des deux suivantes, en vertu de ce qui précède~:
\begin{enumerate}
  \item[(i)] $\mathfrak mB$ \emph{est un idéal de définition de} $B$~;
  \item[(ii)] $B/\mathfrak n$ \emph{est une extension de rang fini du corps}
  $A/\mathfrak m$.
\end{enumerate}
Lorsqu'il en est ainsi, tout $B$-module de type fini est évidemment un
$A$-module quasi-fini.
\end{env}

\begin{corollary}[7.4.5]
\label{0.7.4.5-fr}
Sous les hypothèses de (7.4.2), si $\mathfrak b$ est l'annulateur du
$B$-module $M$, $B/\mathfrak b$ est un $A$-module quasi-fini.
\end{corollary}

Supposons $M\neq0$ (sinon le corollaire est évident). On peut considérer $M$
comme un module sur l'anneau local noethérien $B/\mathfrak b$~; son annulateur
étant alors réduit à $0$, la démonstration de (7.4.2) montre que
$\mathfrak m(B/\mathfrak b)$ est un idéal de définition de
$B/\mathfrak b$. Par ailleurs, $M/\mathfrak nM$ est un espace vectoriel de
rang fini sur $A/\mathfrak m$, étant un quotient de $M/\mathfrak mM$, qui est
par hypothèse de rang fini sur $A/\mathfrak m$~; comme $M\neq0$, on a
$M\neq\mathfrak nM$ en vertu du lemme de Nakayama~; comme
$M/\mathfrak nM$ est un espace vectoriel $\neq0$ sur $B/\mathfrak n$, le fait
qu'il est de rang fini sur $A/\mathfrak m$ implique que $B/\mathfrak n$ est
aussi de rang fini sur $A/\mathfrak m$~; la conclusion résulte donc de (7.4.4)
appliqué à l'anneau $B/\mathfrak b$.

\subsection{Anneaux de séries formelles restreintes}
\label{subsection:0.7.5-fr}

\begin{env}[7.5.1]
\label{0.7.5.1-fr}
Soient $A$ un anneau topologique, linéairement topologisé, séparé et complet~;
soit $(\mathfrak J_\lambda)$ un système fondamental de voisinages de $0$ dans
$A$ formé d'idéaux (ouverts), de sorte que $A$ s'identifie canoniquement à
$\varprojlim A/\mathfrak J_\lambda$ (7.2.1). Pour tout $\lambda$, soit
$B_\lambda=(A/\mathfrak J_\lambda)[T_1,\ldots,T_r]$, où les $T_i$ sont des
indéterminées~; il est clair que les $B_\lambda$ forment un système projectif
d'anneaux discrets. Nous poserons
$\varprojlim B_\lambda=A\{T_1,\ldots,T_r\}$, et nous allons voir que cet anneau
topologique est indépendant du système fondamental d'idéaux
$(\mathfrak J_\lambda)$ considéré. De façon précise, soit $A'$ le sous-anneau
de l'anneau de séries formelles $A[\![T_1,\ldots,T_r]\!]$ formé des séries
formelles $\sum_\alpha c_\alpha T^\alpha$ (avec
$\alpha=(\alpha_1,\ldots,\alpha_r)\in\mathbf{N}^r$) telles que
$\lim_\alpha c_\alpha=0$ (suivant le filtre des complémentaires des parties
finies de $\mathbf{N}^r$)~; nous dirons que ces séries sont les séries
formelles \emph{restreintes} en les $T_i$, à coefficients dans $A$.

